% Editable educator guide. Compile from this directory with:
% pdflatex -interaction=nonstopmode -halt-on-error study-guide.tex
% No external artwork, bibliography processor, or shell escape is required.
\documentclass[11pt,letterpaper]{article}
\usepackage[T1]{fontenc}
\usepackage[utf8]{inputenc}
\usepackage{lmodern}
\usepackage[margin=0.78in,headheight=14pt,headsep=20pt,footskip=25pt]{geometry}
\usepackage{xcolor,amsmath,amssymb,booktabs,tabularx,enumitem,fancyhdr,titlesec}
\usepackage[colorlinks=true,urlcolor=teal,linkcolor=teal]{hyperref}
\definecolor{teal}{HTML}{087E83}
\definecolor{ink}{HTML}{193D42}
\definecolor{muted}{HTML}{58696C}
\renewcommand{\familydefault}{\sfdefault}
\setlength{\parindent}{0pt}
\setlength{\parskip}{7pt}
\renewcommand{\baselinestretch}{1.06}
\setlist[itemize]{leftmargin=16pt,itemsep=4pt,topsep=3pt}
\setlist[enumerate]{leftmargin=19pt,itemsep=6pt,topsep=3pt}
\titleformat{\section}{\Large\bfseries\color{ink}}{}{0pt}{}
\titleformat{\subsection}{\normalsize\bfseries\color{teal}}{}{0pt}{}
\titlespacing*{\section}{0pt}{0pt}{11pt}
\titlespacing*{\subsection}{0pt}{12pt}{3pt}
\pagestyle{fancy}\fancyhf{}
\fancyhead[L]{\small\color{muted}CAST\enspace /\enspace Causal educator kit}
\fancyfoot[L]{\small\color{muted}Synthetic examples \enspace\textperiodcentered\enspace Editable teaching draft}
\fancyfoot[R]{\small\color{muted}\thepage\,/\,6}
\renewcommand{\headrulewidth}{0pt}
\renewcommand{\footrulewidth}{0pt}
\newcommand{\E}{\mathbb{E}}
\newcommand{\Prb}{\mathbb{P}}
\newcommand{\ind}{\mathbf{1}}
\newcommand{\term}[1]{\textbf{\color{ink}#1}}
\newcommand{\source}[1]{{\small\color{muted}#1}}
\pdfstringdefDisableCommands{\def\enspace{ }}
\hypersetup{pdftitle={CAST: Study guide for causal educators},pdfsubject={Potential outcomes, survival effects, causal survival forests and CAST}}
\begin{document}
\thispagestyle{fancy}
{\color{teal}\small STUDY GUIDE \enspace /\enspace 60--90 MINUTES}\par
\vspace{6pt}
{\fontsize{29}{34}\selectfont\bfseries\color{ink}One person,\newline two possible paths}\par
\vspace{5pt}
{\large A lesson on causal effects across time and people}

Hector receives treatment and improves. What would have happened without it?
That missing comparison is the starting point. A survival study then asks what
can be learned about an \emph{average} effect, at a specified time, for a specified
population.

\subsection{Learning goals}
By the end, learners should be able to distinguish a before--after change from
a causal contrast; name a survival estimand and its assumptions; interpret
horizon estimates and a fitted trajectory; and explain why a conditional
average is not an observed individual treatment effect.

\subsection{A suggested class sequence}
\renewcommand{\arraystretch}{1.28}
\begin{tabularx}{\linewidth}{@{}p{0.85in}X@{}}
\toprule
\term{Minutes} & \term{Activity} \\
\midrule
0--10 & Describe Hector's missing outcome. Ask which comparison a before--after plot actually makes. \\
10--25 & Define the survival snapshot and separate individual, sample and population effects. \\
25--40 & Discuss treatment assignment, censoring and overlap. Identify what adjustment must assume. \\
40--55 & Explore the five horizons. Compare the CSF points, CAST curve and simulated answer key. \\
55--70 & Run or inspect the small R lab. Use the frozen results if class time or computing access is limited. \\
70--90 & Optional: critique the fit, discuss heterogeneity and write a causal roadmap for a proposed study. \\
\bottomrule
\end{tabularx}

\subsection{Before class}
Learners need basic probability and a distinction between association and
causation. R is optional. The kit includes editable slides, code, a synthetic
cohort and separate answer key. Its README gives run and compilation commands.

\subsection{The question comes before the method}
Define the population, baseline strategies, outcome and follow-up. State how
the data and assumptions connect to that question. Then estimate and interpret
with uncertainty. In this teaching adaptation of the causal roadmap, CSF and
CAST belong at the estimation step.

\source{Roadmap: Petersen \& van der Laan (2014),
\href{https://doi.org/10.1097/EDE.0000000000000078}{doi:10.1097/EDE.0000000000000078}.
Implementation authority: \href{https://github.com/ishuryak/cast_demo_website}{Igor Shuryak's demo repository}.}

\newpage
\section{1\enspace The missing outcome and the snapshot}
\term{Potential outcomes} describe what would occur under specified alternative
treatment choices. Let $W_i\in\{0,1\}$ be person $i$'s baseline treatment and
$T_i(w)$ their potential event time under choice $w$. At horizon $t$, define
\[
 A_i^w(t)=\ind\{T_i(w)>t\},\qquad
 \delta_i(t)=A_i^1(t)-A_i^0(t).
\]
$A_i^w(t)$ is a survival indicator: 1 means alive beyond $t$, 0 means the event
has occurred by $t$. The individual effect $\delta_i(t)$ is $-1$, 0 or $1$.
We cannot observe both potential outcomes for the same person at the same
horizon. Randomization does not fill in that missing value.

Under \term{consistency}, the event time on the assigned path is
$T_i=T_i(W_i)$. An observed improvement before versus after treatment compares
different times on that one path. It does not compare treatment choices at
the same time.

\subsection{Three averages that should not be conflated}
\renewcommand{\arraystretch}{1.48}
\begin{tabularx}{\linewidth}{@{}p{1.24in}X@{}}
\toprule
\term{Realized sample} &
$\displaystyle\Delta_{\mathrm{SATE}}(t)=\frac1n\sum_{i=1}^n
 [A_i^1(t)-A_i^0(t)]$\newline
The average of the realized, paired potential outcomes in this sample; some are unobserved. \\
\term{Population} &
$\displaystyle\Delta(t)=\Prb\{T(1)>t\}-\Prb\{T(0)>t\}$\newline
The survival-probability difference in a stated target population. \\
\term{Conditional} &
$\displaystyle\tau(t,x)=\E[A^1(t)-A^0(t)\mid X=x]$\newline
The average contrast for people with baseline profile $x$; it is not $\delta_i(t)$. \\
\bottomrule
\end{tabularx}

\subsection{What the simulation's answer key means}
The generator supplies conditional survival probabilities $S_w(t\mid x)$.
The answer key averages their contrast over the generated baseline profiles:
\[
 \Delta_{\mathrm{oracle},X}(t)=\frac1n\sum_{i=1}^n
 \{S_1(t\mid X_i)-S_0(t\mid X_i)\}.
\]
This is a sample-standardized \emph{probability} contrast, also interpretable
as the expected SATE given these profiles under the generating model. It is
not the realized SATE and does not reveal whose outcome treatment changed.
For simulations with a hidden baseline factor, the oracle also uses that
factor; the fitted models do not. The supplied small cohort has no hidden
confounding.

\subsection{Read the units before reading the direction}
A difference of $0.10$ is \term{10 percentage points} (pp). It is neither a
10\% relative increase nor a hazard ratio. Positive values favor treatment
on this survival scale. Hector's opening sketch uses symptom burden, where
lower is better; that vertical direction does not carry over to survival.

\source{Potential outcomes and average effects:
Hern\'an \& Robins, \href{https://miguelhernan.org/whatifbook}{Causal Inference: What If}, chapters 1--3.}

\newpage
\section{2\enspace What makes a comparison causal?}
\term{Confounding} occurs when treatment groups differ in causes of the
outcome. In this simulation, baseline health affects both treatment assignment
and survival. Comparing the observed groups without adjustment can therefore
mix treatment effects with pre-existing differences.

\subsection{Observed survival data}
Let $C_i$ be a censoring time. The recorded data include baseline covariates
$X_i$, treatment $W_i$, observed follow-up $Y_i$ and event indicator $D_i$:
\[
 Y_i=\min\{T_i,C_i\},\qquad D_i=\ind\{T_i\le C_i\}.
\]
If $D_i=0$ and $Y_i<t$, survival at $t$ is unknown even on the assigned path.
Censoring and the missing counterfactual are different sources of missing
information. Censored people are not automatically events or survivors at
every later horizon.

\subsection{Identification assumptions to discuss}
\begin{description}[leftmargin=0pt,labelsep=0pt,itemsep=8pt,style=nextline]
\item[\term{Well-defined strategies and consistency}]
Treatment has a specified meaning; the observed outcome follows the assigned
strategy. Interference between people must be absent or explicitly modeled.
\item[\term{Conditional exchangeability}]
Given adequate baseline covariates $X$, treatment assignment carries no
remaining information about potential event times:
$\{T(0),T(1)\}\perp W\mid X$.
Measuring many variables does not establish this assumption.
\item[\term{Treatment positivity / overlap}]
For relevant profiles, both treatment choices are possible:
$0<e(x)<1$, where $e(x)=\Prb(W=1\mid X=x)$ is the
\term{propensity score}. Sparse overlap makes adjustment unstable.
\item[\term{Censoring assumptions and follow-up support}]
A sufficient working assumption is $T\perp C\mid X,W$, with positive
probability of remaining uncensored through the horizons of interest.
Relevant common causes of censoring and the event must be handled.
\end{description}

\subsection{Identification is separate from estimation}
These assumptions connect a causal target to the distribution of observed
data. An estimator then learns that distribution with finite data and modeling
error. Flexible forests cannot repair unmeasured confounding or create support
where a treatment is never observed. Generalizing to a different population
also needs a defensible sampling or transport argument.

\term{Challenge prompt.} A model fits the observed data very well and its
confidence interval is narrow. What evidence would still be needed before
calling the estimate causal? Name one plausible unmeasured common cause and
one overlap check for your proposed study.

\source{CSF assumptions: Cui et al. (2023),
\href{https://doi.org/10.1093/jrsssb/qkac001}{doi:10.1093/jrsssb/qkac001}.
Software target: \href{https://grf-labs.github.io/grf/reference/causal_survival_forest.html}{grf causal\_survival\_forest reference}.}

\newpage
\section{3\enspace Effects across time and people}
\term{Causal survival forests (CSF)} estimate conditional average effects with
right-censored outcomes. This demo uses the survival-probability target,
evaluated at five horizons, then reports a cohort average at each horizon.
A \term{horizon} is a chosen evaluation time; it is not a bin for event times.

\subsection{What CAST adds}
CAST (Causal Analysis for Survival Trajectories) models how the horizon
estimates vary over follow-up. In this demo, a quadratic is fitted to the CSF
estimates:
\[
 \widehat\Delta_{\mathrm{CAST}}(t)=\widehat\beta_0+
 \widehat\beta_1t+\widehat\beta_2t^2.
\]
More points drawn along this curve are predictions from that fit, not new
observations or additional forest estimates. Keep the original horizon points
visible: smoothing can miss a peak, a crossing or another feature.

\subsection{Optional math: correlated estimates and the band}
Let $\widehat{\boldsymbol d}$ contain the horizon estimates, $B$ have rows
$b(t_j)^\top=(1,t_j,t_j^2)$, and $\widehat\Sigma$ estimate their joint
covariance. With fitting weights $Q$,
\begin{align*}
 \widehat{\boldsymbol\beta}
   &=(B^\top QB)^{-1}B^\top Q\widehat{\boldsymbol d},\\
 \widehat V_\beta
   &=(B^\top QB)^{-1}B^\top Q\widehat\Sigma QB(B^\top QB)^{-1}.
\end{align*}
The shared people induce dependence across horizons. This implementation
forms $\widehat\Sigma$ from patient-aligned influence scores and applies
covariance shrinkage. It uses generalized least squares (GLS),
$Q=\widehat\Sigma^{-1}$, when conditioning and fit checks pass; otherwise it
uses weighted least squares (WLS) with inverse-variance diagonal weights.
A reported 95\% pointwise band at $t$ is
\[
 b(t)^\top\widehat{\boldsymbol\beta}\ \pm\
 1.96\sqrt{b(t)^\top\widehat V_\beta b(t)}.
\]
\term{Pointwise} does not mean 95\% coverage of the whole trajectory at once.
The calculation also does not account for all uncertainty from choosing the
fit or from a misspecified quadratic. The full procedure's coverage needs
separate validation. These implementation choices belong to the demo code;
they are not all specified in the original CAST preprint.

\subsection{Heterogeneity and the meaning of time}
\term{Treatment-effect heterogeneity} means $\tau(t,x)$ varies across profiles
$x$, on a stated outcome scale and at a stated horizon. Averaging over a
target population gives $\Delta(t)=\E_X[\tau(t,X)]$. The website plots cohort
averages; it does not display a subgroup analysis or identify Hector's effect.

A falling positive survival difference still favors treatment at that horizon.
Following the effect of a \emph{baseline} treatment is different from deciding
when to switch or stop it. A fitted peak is not a treatment rule.

\source{CAST framework: \href{https://arxiv.org/abs/2505.06367}{Yang et al. (2025)}.
Exact covariance, weighting and band:
\href{https://github.com/ishuryak/cast_demo_website/blob/43a00101100801eb6c048800d24788455fac0f8b/R/cast_core.R}{versioned R/cast\_core.R}.}

\newpage
\section{4\enspace Class exercises}
\subsection{A. A contrast is not a count of identifiable beneficiaries}
In a hypothetical population of 10 people, suppose the average survival
effect at one horizon is $+0.20$. Let $H$ count people with
$\delta_i(t)=+1$ and $L$ count people with $\delta_i(t)=-1$.
\begin{enumerate}
\item Write the relationship among $H$, $L$ and the average effect.
\item Give two different pairs $(H,L)$ that produce $+0.20$.
\item Does this average tell you that a particular two people were saved?
Explain what remains unobserved.
\end{enumerate}

\subsection{B. Read a frozen website example}
These values come from the 2,000-person \emph{reversal} scenario with measured
confounding strength 1 and hidden-confounding strength 0. All entries are
survival-probability differences in pp, rounded from the frozen export.
\renewcommand{\arraystretch}{1.28}
\begin{center}
\begin{tabular}{@{}rrrrr@{}}
\toprule
\term{Month} & \term{Unadjusted} & \term{CSF} & \term{CAST} & \term{Oracle} \\
\midrule
12  & 13.49 & 6.51 & 6.39 & 7.51 \\
36  & 32.35 & 18.09 & 10.34 & 16.38 \\
60  & 28.56 & 11.42 & 9.96 & 9.15 \\
84  & 14.51 & $-0.12$ & 5.28 & $-1.59$ \\
108 & 6.81 & $-4.38$ & $-3.72$ & $-5.60$ \\
\bottomrule
\end{tabular}
\end{center}
\begin{enumerate}[resume]
\item At 108 months, which comparisons favor treatment? Interpret the CSF
estimate in words. Does its sign tell you when treatment should be stopped?
\item At which of the five horizons is the absolute CAST--oracle gap largest?
Compute it. What does this example show about smoothing?
\end{enumerate}

\subsection{C. Inspect or run the supplied cohort}
The kit's \texttt{data/demo-cohort.csv} contains \term{600 synthetic people},
generated with seed 20260905, reversal effects, measured confounding strength
1 and no hidden confounding. It is a distinct sample from the frozen website
export. Baseline columns are \texttt{age, stage, ps, comorb, smoke, sex,
ethnicity}; treatment is \texttt{W}, follow-up in months is \texttt{Y}, and
the event indicator is \texttt{D}. Read category coding in the kit README.

Run \texttt{class-lab.R} from the extracted kit as described in its README.
It fits Igor's bundled implementation to the cohort. Compare the five horizon
estimates with \texttt{data/answer-key.csv}, which contains the generating
model's probability contrasts for these same profiles. An estimate is not
expected to equal the answer key in one finite sample. Intervals can miss it
too: see the kit README.

\term{Discuss.} Describe one baseline imbalance and the fraction censored.
Choose a clinically meaningful horizon \emph{before} examining its result.
Then state the population, comparison, assumptions and uncertainty in a
two-sentence interpretation. Use \texttt{data/frozen-horizons.csv} for a
no-compute comparison across the website's 24 scenarios.

\newpage
\section{5\enspace Answer notes and further reading}
\subsection{Answers to the numbered exercises}
\begin{enumerate}
\item The realized average is $(H-L)/10$, so $H-L=2$.
\item $(2,0)$ and $(4,2)$ both give $+0.20$, with different numbers helped
and harmed.
\item No. Even an exact average would not identify the missing potential
outcomes. An estimated probability difference also has sampling and modeling
uncertainty.
\item At 108 months the unadjusted comparison is positive; CSF, CAST and the
oracle are negative. CSF estimates survival under baseline treatment to be
4.38 pp lower than under control for the cohort target. This is a horizon
contrast, not a comparison of continuing versus stopping treatment then.
\item At 84 months: $|5.28-(-1.59)|=6.87$ pp. The gap at 36 months is
6.04 pp. A smooth fit can miss the generating trajectory even if it looks
plausible; it need not improve on the separate horizon estimates.
\end{enumerate}

\subsection{Facilitator notes for the lab and discussion}
The cohort has one row per person, not paired potential outcomes. The answer
key averages known model probabilities; do not use it as fitting data.
The fraction with $D=0$ describes recorded censoring, not whether survival is
known at every horizon.

An interpretation should state the baseline choices, target profiles, horizon,
units and uncertainty. Separate identification failures, such as hidden
confounding, from estimation problems, such as poor fit. For a new population,
revisit the averaging distribution. For changing treatment over time, define
a new longitudinal question and revisit identification.

\subsection{A compact vocabulary recap}
\term{Estimand}: the quantity we want to learn. \term{Identification}: relating
that quantity to observed-data distributions under assumptions.
\term{Estimator}: the procedure used to learn it from a sample.
\term{Counterfactual}: an outcome under an alternative specified treatment.
\term{Censoring}: incomplete event-time observation.
\term{Overlap}: support for both treatment choices within relevant profiles.
\term{CSF}: a forest for conditional average effects with censored outcomes.
\term{CAST}: a model for an effect trajectory across horizons.
\term{Pointwise interval}: uncertainty at one specified time.
\term{Heterogeneity}: variation in a conditional average effect across profiles.

\subsection{Sources and editable materials}
{\small
\begin{itemize}[itemsep=3pt]
\item Hern\'an \& Robins (2020). \href{https://miguelhernan.org/whatifbook}{\emph{Causal Inference: What If}}. Potential outcomes, identification and generalization.
\item Cui et al. (2023). \href{https://doi.org/10.1093/jrsssb/qkac001}{\emph{Estimating heterogeneous treatment effects with right-censored data via causal survival forests}}. JRSS B 85:179--211.
\item Yang et al. (2025). \href{https://arxiv.org/abs/2505.06367}{\emph{CAST: Time-Varying Treatment Effects}}. Methods preprint, arXiv:2505.06367.
\item Petersen \& van der Laan (2014). \href{https://doi.org/10.1097/EDE.0000000000000078}{\emph{Causal models and learning from data}}. Epidemiology 25:418--426.
\item \href{https://grf-labs.github.io/grf/reference/causal_survival_forest.html}{grf CSF reference} and \href{https://github.com/ishuryak/cast_demo_website}{Igor Shuryak's demo repository}: target definitions, source code and implementation details.
\end{itemize}}
\end{document}
